\section{Optimization Perspective on Lewis Weights}
\label{sec:convex}

Although simple and appealing, the iterative scheme described fails
when $p \geq 4$, since $\textsc{LewisIterate}$ stops being a contraction (further,
runtime approaches infinity as $p \to 4$).  To obtain Lewis weights for general
$p$, we therefore need another approach.  Here, we will give a characterization of
Lewis weights based on the solution to an optimization problem.

This optimization problem is essentially derived from a standard
proof of the existence of Lewis weights (\cite{BanachBook}, III.B, 7).  In fact,
it is interesting to note that both of these algorithms correspond fairly
directly to proofs of the existence and uniqueness of Lewis weights.
The optimization argument given here is valid for
all $p$, but only directly leads to an efficient algorithm for $p \geq 2$
(this is unimportant, since the iterative scheme is extremely efficient
when $p < 2$).

The optimization problem we will look at is, over symmetric matrices $\MM$,
\begin{equation*}
\begin{aligned}
& \underset{\MM}{\text{maximize}} & & \det M \\
& \text{subject to}
& & \sum_i (\aa_i^T \MM \aa_i)^{p/2} \leq d, \\
&&& \MM \succeq 0.
\end{aligned}
\end{equation*}

Note that when $p \geq 2$, the region in question is convex: it is the intersection
of the (convex) positive semidefinite cone and and $\ell_{p/2}$-norm constraint on
a vector of linear functions of $\MM$ (the $\aa_i^T \MM \aa_i$).  Maximizing the
determinant is also equivalent to minimizing the convex function $-\log \det \MM$,
so it is a convex optimization problem (we will give a specific algorithm using
convex optimization tools to find approximate Lewis weights below, in Theorem~\ref{thm:polyslow}.

\begin{lemma}
The optimization problem given above attains its maximum.
Further, for any matrix $\QQ$ reaching this maximum, the weights
\begin{equation*}
\ww_i = (\aa_i^T \QQ \aa_i)^{p/2}
\end{equation*}
satisfy the definition (Definition~\ref{def:lewis}) of Lewis weights.
\end{lemma}
\begin{proof}
First, note that the region in question is compact.  Since the function to be
optimized is continuous, it must attain a maximum.

Now, consider any matrix $\QQ$ that attains the maximum.
Then $\QQ$ will saturate the $\sum_i (a_i^T \MM a_i)^{p/2}$ constraint since otherwise
it could just be scaled up, but will be positive definite (not on the boundary of
the semidefinite cone) because otherwise it would have determinant 0.  We may then
say that $\QQ$ is a local maximum of the determinant subject to the constraint that
\begin{equation*}
\sum_i (\aa_i^T \MM \aa_i)^{p/2} = d.
\end{equation*}
Since all the functions in question are smooth, we may characterize such a local optimum
by Lagrange multipliers.  The gradient of the constraint (note that we are in the vector
space of symmetric matrices, so this is a matrix) at $\QQ$ can be seen to be
\begin{equation*}
p/2 \sum_i (\aa_i^T \QQ \aa_i)^{p/2 - 1} \aa_i \aa_i^T.
\end{equation*}
The gradient of the determinant, on the other hand, at $\QQ$, is $\det(\QQ) \QQ^{-1}$.
Lagrange multipliers then imply that at $\QQ$, $\QQ^{-1}$ must be parallel to
$\sum_i (\aa_i^T \QQ \aa_i)^{p/2 - 1} \aa_i \aa_i^T$.

The claim is that
\begin{equation*}
\ww_i = (\aa_i^T \QQ \aa_i)^{p/2}.
\end{equation*}
then must satisfy the conditions for Lewis weights.  To argue this, we note that
for these $\ww_i$, $\AA \WW^{1 - 2/p} \AA$ is
\begin{equation*}
\sum_i \ww_i^{1-2/p} \aa_i \aa_i^T = \sum_i (\aa_i^T \QQ \aa_i)^{p/2 - 1} \aa_i \aa_i^T.
\end{equation*}
In other words, $\AA \WW^{1 - 2/p} \AA$ is equal to $C \QQ^{-1}$ for some scalar $C$.
That implies that the $\ww_i$ could be defined as
\begin{equation*}
C^{p/2} (\aa_i^T (\AA \WW^{1 - 2/p} \AA)^{-1} \aa_i)^{p/2}.
\end{equation*}
Then some scaling of the $\ww_i$ satisfies
$\ww_i = (\aa_i^T (\AA \WW^{1 - 2/p} \AA)^{-1} \aa_i)^{p/2}$,
so $\ww$ is a multiple of Lewis weights. But since
\begin{equation*}
\sum_i \ww_i = \sum_i (\aa_i^T \MM \aa_i)^{p/2} = d
\end{equation*}
$\ww_i$ defined this way must be precisely the Lewis weights.
\end{proof}

\begin{corollary}
\label{cor:conlewis}
For all $p$, there exists an assignment of weights satisfying Definition~\ref{def:lewis}.
Furthermore, for $p \geq 2$, this assignment is unique.
\end{corollary}
\begin{proof}
The existence statement just follows from taking any of the $\QQ$ attaining the maximum.
For $p \geq 2$, the uniqueness follows from the fact that a strictly convex function
(which $-\log \det \MM$ is) attains a unique minimum on a convex set, and that applying
the argument in reverse every set of weights satisfying Definition~\ref{def:lewis}
is induced by such a $\QQ$.
\end{proof}
Of course, as mentioned earlier, uniqueness for $p < 2$ follows from Corollary~\ref{cor:itlewis}.

It is worth noting that this optimization problem gives a simple, geometric
characterization of the Lewis weights.  $\sum_i (\aa_i^T \MM \aa_i)^{p/2}$
is proportional to the average value of $\norm{\AA x}_p^p$ inside the ellipsoid
$x^T \MM^{-1} x \leq 1$.  Thus, the quadratic form induced by the Lewis weights
corresponds to the ellipsoid of maximum volume with a limited $p$-moment of
$\norm{\AA x}_p$; it is a ``softer'' version of a John ellipsoid, where the max
(effectively the $\infty$-moment) of $\norm{Ax}_p$ would be limited instead.

To use this algorithmically, we only assume an approximate solution.
We can analyze such a solution using the following lemma:
\begin{lemma}
There exists a constant $C$ such that for any $0 < \epsilon \leq 1$,
and positive semidefinite matrix $\MM$ satisfying
$\sum_i (\aa_i^T \MM \aa_i)^{p/2} = d$ and $\det \MM \geq (1-C \epsilon^2) \det \QQ$,
$\MM \approx_{1+\epsilon} \QQ$.
\end{lemma}
\begin{proof}
First note that $\QQ$ must be a maximum of
$\trace{\MM \QQ^{-1}}$ within the feasible region (otherwise, there would be a
direction in which the determinant increases around $\QQ$, making it not a local
maximum).  But for any matrix $\MM$ with $\trace{\MM \QQ^{-1}} \leq d$, if
$\MM \QQ^{-1}$ has any eigenvalues further than $\epsilon$ from 1,
$\det(\MM^{-1}) \leq (1 - O(\epsilon^2)) \det(\QQ^{-1})$.  Therefore,
an $O(\epsilon^2)$-approximate solution to the optimization problem implies
an $\epsilon$-approximation of $\QQ$ and of the Lewis weights themselves.
\end{proof}

The algorithmic guarantee is then:
\begin{theorem}
\label{thm:polyslow}
There exists a function $f(p, \epsilon)$ and a constant $C$ such that for any $\AA$, $p \geq 2$,
$(1+\epsilon)$-approximate Lewis weights (or the quadratic form $\QQ$) can be computed in time $O(f(p, \epsilon) n \log(n) \log \log(n) d^C)$.
\end{theorem}
Note that the $p < 2$ case already follows from Lemma~\ref{lem:computeLewis}.
\begin{proof}
For any positive real number $D$, consider the intersection of the sets $\MM \succeq 0$,
$\sum_i (\aa_i^T \MM \aa_i)^{p/2} \leq d$, and $\det(\MM) \geq D$.  This is a convex set
with a polynomial time separation oracle.

Furthermore, (if $D \leq \det \QQ$) this set always contains $\QQ$, which satisfies
$\QQ \preceq n^{1-2/p} \AA^T \AA$ by the argument from Lemma~\ref{lem:initialization}
(which nowhere assumes that $p < 4$) and is therefore contained in the ellipsoid
(over matrices!) $\norm{\QQ (\AA^T \AA)^{-1}}_F^2 \leq n^{2-4/p} d$.  If $D$ is within a
constant factor of $\det \QQ$, then the entire intersection is contained in a constant multiple
of that ellipsoid, and the volume ratio between it and the ellipsoid is at most $n^{O(d^2)}$.
Thus, for any such $D$, the ellipsoid algorithm can find an element in time $d^C \log n$ iterations,
each of which will take $n d^C$ time.

Finally, one may binary search (over an exponentially spaced set of possible $D$ values) to find
such a $D$ within an appropriate constant factor $1+O(\epsilon^2)$ in $\log \log n + \log d$ iterations.
\end{proof}
This is a polynomial-time algorithm, but not an input-sparsity one.
